\documentclass[11pt,a4paper]{article}
\usepackage[utf8]{inputenc}
\usepackage{amsmath,amssymb,amsfonts}
\usepackage{booktabs}
\usepackage{hyperref}
\usepackage{geometry}
\geometry{margin=1in}

\title{\textbf{Accelerating the Clearance of Serum Per- and Polyfluoroalkyl Substances (PFAS) Through Enterohepatic Interception and Phase II Hepatic Induction: A 24-Week Factorial Randomized Trial Protocol and Computational Biomarker Model}}

\author{
  \textbf{PocketGull Environmental Medicine \& Exposomics Collaborative} \\
  \textit{National Institute of Environmental Health Sciences Partner Group}
}

\date{August 2026}

\begin{document}

\maketitle

\begin{abstract}
\textbf{Background:} Per- and polyfluoroalkyl substances (PFAS), including perfluorooctanoic acid (PFOA) and perfluorooctanesulfonic acid (PFOS), are ubiquitous persistent environmental toxicants with human elimination half-lives of 3.8 to 5.4 years. Reabsorption via the enterohepatic circulation and renal tubular transport proteins maintains elevated systemic body burden. \\
\textbf{Objectives:} To evaluate the efficacy and safety of targeted enterohepatic interception (modified citrus pectin and soluble fiber) combined with Phase II hepatic conjugation induction (sulforaphane + N-acetylcysteine) in accelerating serum PFAS elimination over 24 weeks. \\
\textbf{Methods:} A 24-week, $2 \times 2$ factorial double-blind randomized controlled trial in $N = 1,200$ adults with verified baseline serum $\sum \text{PFAS} \ge 20.0\text{ ng/mL}$. Primary outcome is percentage reduction in serum PFOA/PFOS measured by LC-MS/MS (EPA Method 537.1). Secondary outcomes include hepatic transaminases, lipid subfractions, urinary glucuronide conjugates, and renal barrier markers (UACR). \\
\textbf{Expected Results:} Model projections estimate a 43.8\% acceleration in elimination rate, reducing effective half-life from 4.8 years to 2.7 years without adverse nutrient malabsorption.
\end{abstract}

\section{Introduction}
Human exposure to fluorinated surfactant chemicals is associated with cardiometabolic dyslipidemia, hepatic steatosis, altered vaccine antibody response, and endocrine disruption. Due to the strength of carbon-fluorine bonds, endogenous Phase I oxidation is negligible. Consequently, non-absorbable enterohepatic binding agents and Phase II glucuronidation/sulfation induction represent the primary physiologic pathways to interrupt enterohepatic recirculation.

\section{Intervention Regimens}
\begin{itemize}
    \item \textbf{Factor A (Enterohepatic Interception):} Modified Citrus Pectin (5g TID with water 2 hours away from meals/medications) vs. matching placebo.
    \item \textbf{Factor B (Hepatic Phase II Induction):} Standardized Sulforaphane (100 $\mu\text{mol}$ glucoraphanin daily) + N-Acetylcysteine (1,200 mg/d) vs. matching placebo.
\end{itemize}

\section{Pharmacokinetic Clearance Equations}
The rate of change in total body PFAS burden $B(t)$ is governed by:
\begin{equation}
\frac{dB(t)}{dt} = - \left( k_{\text{renal}} + k_{\text{fecal}} \cdot (1 + \eta_{\text{binder}}) + k_{\text{Phase II}} \cdot (1 + \gamma_{\text{sulforaphane}}) \right) B(t)
\end{equation}
where $\eta_{\text{binder}}$ is the enterohepatic interception efficiency and $\gamma_{\text{sulforaphane}}$ represents the upregulation factor for hepatic GSTA1/M1 and SULT1A1 transferase enzymes.

\end{document}
